\chapter{Chapter 7 --- Conclusion}\label{chapter-7-conclusion}

This thesis investigated how to accelerate test-time discovery in complex code generation via execution-guided self-distillation. The primary methodological contribution is a teacher-first organization: the self-teacher generates trajectories under the guidance of execution feedback, after which a verifier and a semantic judge filter them for functional correctness and structural independence. By explicitly removing the fallback mechanism, the framework ensures a copied solution is never distilled. This positions the proposed method at the critical intersection between on-policy Self-Distillation Policy Optimization (SDPO) and off-policy Supervised Fine-Tuning (SFT).

The empirical evaluation presents a robust and compute-fair performance profile. On a medium-scale matched evaluation, teacher-first significantly outperforms the student-first baseline and successfully escapes the flat-reward trap on the hardest programming problems. Under a matched generation budget and across a compute-to-correct frontier analysis, this performance advantage holds consistently, proving that the gain is an algorithmic property of the chosen distillation trajectory rather than a sampling-volume artifact. Furthermore, the investigation reveals that the reprompt-template formulation of the execution feedback heavily influences discovery, exhibiting its most pronounced effects on hard problems. 

With reasoning traces enabled, distillation from a feedback-conditioned context measurably suppresses epistemic verbalization in code generation, with the structural implications of this suppression proven to be strictly regime-dependent. Most decisively, by holding the model architecture fixed and supplying a method-bearing reference, the model achieves escape on mathematical tasks where an answer-only reference fails. This explicitly validates the central value-versus-procedure boundary: successful test-time escape requires the privileged reference to encode an in-reach method rather than a superficial final value.

While this work remains an empirical characterization conducted within a single model family at a personal compute scale—where scaling to stronger architectures is expected to naturally narrow the baseline margin—the consolidated contribution provides a compute-fair, mechanistically explained, and boundary-validated account of execution-guided self-distillation. The framework successfully mapping out when and how the approach accelerates test-time discovery provides a comprehensive foundation for formal publication and future extensions in inference-time optimization.